\documentclass[10pt]{article}
\usepackage[margin=20mm]{geometry}
\usepackage{amsmath,amssymb,mathtools,fontspec,unicode-math,xurl,hyperref}
\IfFontExistsTF{DejaVu Serif}{\setmainfont{DejaVu Serif}}{\setmainfont{Latin Modern Roman}}
\setmathfont{latinmodern-math.otf}
\setlength{\parindent}{0pt}\setlength{\parskip}{6pt}\sloppy
\title{Extensions, inverse regularity, and jet formulas in the theory of half-Lie groups}
\author{AI-assisted technical review}
\date{}
\begin{document}
\maketitle
Abstract: The complete 42-page author manuscript of Bauer, Harms, and Michor, arXiv:2302.01631v3, is examined. Explicit counterexamples invalidate the joint-smoothness conclusion of Theorem 5.10(a) and the descent assertion of Proposition 5.12. The proof of the no-loss theorem uses a false inverse formula; a short-time repair is available, but its extension to the full interval is not established here. A further Major gap is the invalid boundedness inference used to construct the Banach model for differentiable elements; a fixed-order repair under the local-addition hypothesis is given. Several local jet, sign, and extension formulas require correction. The strong-metric Hopf–Rinow argument on the original group is assessed with its stated additional hypotheses intact. No counterexample to that result, or to the main differentiable- element theorem, is asserted. Scientific adjudication remains pending.\par
\section*{Identity, scope, and dependencies}
The reviewed work is Martin Bauer, Philipp Harms, and Peter W. Michor, Regularity and completeness of half-Lie groups, DOI 10.4171/JEMS/1587. The primary scientific source is the author preprint arXiv:2302.01631v3, dated 7 January 2025 [paper]. During the initial review, its 42 pages, including all proofs and references, were read, and all 42 rendered source pages were actually viewed. Page references below are one-based PDF pages and agree with the printed numbering.\par
The journal page records acceptance on 4 December 2023 and publication on 10 January 2025 [journal]. The author’s comment that v3 agrees with the published version is not an independent comparison. The subscription journal PDF was not acquired for a bytewise or pagewise comparison. The retained v3 TeX was read selectively to resolve formulas; it was not compiled or executed. No separate article-specific supplement or computational code was found in the bounded acquired-source search. Their absence is not a defect.\par
A right half-Lie group has continuous group operations and smooth individual right translations; joint smooth multiplication is not assumed (Definition 2.1, p. 3). Definition 3.1 (pp. 4–5) requires both \(L_x\) and \(L_{x^{-1}}\) to be \(C^k\). Theorem 3.4 and Lemma 3.7 construct \(G^k\) and the ILB Lie group \(G^\infty\) under a right-invariant local addition; their proofs run through Sections 9–10. Theorem 4.2 uses these spaces to integrate smooth controls. Definition 4.1 requires existence and uniqueness of evolution, not smooth dependence of the evolution map on the control. No extra Milnor-regularity condition is imposed.\par
Section 5 develops extension data and differentiable elements in extensions; Section 6 gives bundle examples. Theorem 7.7(a)–(d) concerns a connected group with a strong, smooth, right-invariant metric. Its geodesic-convexity conclusion (e) additionally assumes \(L^2\)-regularity and weak closedness of every endpoint fibre \(\mathcal A_x\). Its final assertion for the weaker metrics on \(G^k\) depends on the no-loss result 7.5–7.6. These dependencies are distinct.\par
\section*{Substantive extension and proof issues}
Problem 1 (Major: individual conjugation does not imply a jointly smooth action).\par
Location. Definitions 5.1–5.2 (pp. 8–9), Remark 5.8 (p. 11), Definition 5.9(b) and Theorem 5.10(a) (p. 12).\par
Let \(H=\ell^2(\mathbb C)\), viewed as a real Hilbert space, and set \[U_t(z_j)_{j\ge1}=(e^{ijt}z_j)_{j\ge1}.\] This is a strongly continuous unitary group. On \(E=\mathbb R\times H\) put \[(t,x)(s,y)=(t+s,U_{-s}x+y),\qquad
 (t,x)^{-1}=(-t,-U_tx).\] The group operations are continuous. For fixed \((s,y)\), right translation is a bounded affine smooth diffeomorphism. Thus \(E\) is a Banach right half-Lie group. Take the additive Banach Lie group \(N=H\), the quotient \(G=\mathbb R\), and \[i(n)=(0,n),\quad p(t,x)=t,\quad s(t)=(t,0),\quad r(t,x)=x.\] The split inclusion and projection have the required initial/final properties, and \(s,r\) are global smooth section and retraction. This is a smooth split extension, stronger than the assumptions of Theorem 5.10(a). For each fixed \(g=(t,x)\), \[g\,i(n)\,g^{-1}=i(U_tn)\] is a smooth bounded-linear automorphism of \(N\), as required by Definition 5.2. The flat local addition on the product vector space is right invariant, since right translations are affine.\par
Nevertheless the canonical section induces \[\alpha_t(n)=U_{-t}n,\qquad f(t,s)=0.\] Definition 5.9(b) requires \(\alpha:V\times N\to N\) to be jointly smooth on an identity neighbourhood \(V\). Take \(n_0=(1/j)_{j\ge1}\in H\). If \(t\mapsto U_{-t}n_0\) were differentiable at zero, its derivative would have every coordinate equal to \(-i\). This sequence is not in \(\ell^2\). Hence the orbit is not even \(C^1\). Since \(N\) is abelian, changing the section only introduces trivial inner conjugations and cannot remove the obstruction.\par
Impact and repair boundary. This is a counterexample to the actual implication in 5.10(a). Individual smooth conjugations, even with a smooth split extension and a right-invariant local addition, do not supply the required joint smoothness. One can restrict the classification to extensions with jointly smooth induced \(\alpha\), or weaken the datum and redo the atlas construction. The latter repair is not established here. The example does not refute the differentiable-vector theorem or Hopf–Rinow: its product flat strong metric is right invariant and complete.\par
Problem 2 (Major: a local addition respecting the kernel need not descend).\par
Location. Proposition 5.12 and its proof (pp. 14–15); the last quotient-local-addition step of Theorem 5.11 (p. 14).\par
Consider the split Lie-group extension \[0\longrightarrow\mathbb R\times\{0\}\longrightarrow(\mathbb R^2,+)
 \xrightarrow{p}(\mathbb R,+)\longrightarrow0,\qquad p(a,b)=b.\] On \(V=\{(a,b;u,v):|u|<1/2,\ |v|<1\}\) define \[\tau_E(a,b;u,v)=\bigl(a+u,\ b+(1+u)v\bigr).\] This is a smooth right-invariant local addition. It fixes the zero section, its domain is invariant under tangent translations, and \((\pi,\tau_E)\) has the smooth inverse \[u=A-a,\qquad v=\frac{B-b}{1+A-a}\] on its open image. It maps \(TN\) to \(N\) and respects the tangent directions of every kernel coset. It is already normalized as in Remark 8.2: its vertical derivative at zero is the identity.\par
But \[p\tau_E(a,b;u,v)=b+(1+u)v,\qquad Tp(a,b;u,v)=(b,v).\] For fixed nonzero \(v\), varying \(u\) changes the first expression without changing the second. At \(v=1/2\), \(u=0\) and \(u=1/4\) give \(b+1/2\) and \(b+5/8\). The obstruction occurs for arbitrarily small nonzero \(u,v\). Therefore no \(\tau_G\) satisfies \(p\tau_E=\tau_G Tp\), even after shrinking to an open neighbourhood of the zero section.\par
Impact and repair boundary. The proof assumes separation of kernel and quotient directions that does not follow from respecting \(N\). Constancy of \(p\tau_E\) on the fibres of \(Tp\), or a separate compatible connection/local addition, is needed. Being totally geodesic is insufficient. This refutes descent of the given addition, not existence of another compatible one: the standard flat addition works here. The quotient step in 5.11 needs its own justification; its entire differentiable-element conclusion is not disproved.\par
Problem 3 (Major: the inverse step in the global no-loss argument is invalid).\par
Location. Theorem 7.5, especially its final proof lines (pp. 23–24); Definition 3.1 (pp. 4–5); Corollary 7.6 and the final assertion of 7.7 (pp. 24–26).\par
Writing \(\gamma(t)=\pi_G\operatorname{Fl}_t^S(X)\), the proof identifies the inverse of \(L_{\gamma(t)}\) with \(L_{\gamma(-t)}\). A second-order flow need not give a one-parameter subgroup. Even a right-invariant spray on \(G=(\mathbb R,+)\) supplies a counterexample: \[S(x,v)=(v,v^2),\quad X=(0,1),\quad
 \gamma(t)=-\log(1-t),\quad
 \gamma(1/2)+\gamma(-1/2)=\log(4/3)\ne0.\] This field is smooth, translation invariant, second order, and quadratic in velocity. The two translations are not inverses. Here \(G^k=G\), so this refutes the formula only, not the no-loss conclusion. The gap matters because Definition 3.1 explicitly requires inverse regularity and does not know it to follow from forward regularity alone.\par
A short-time repair is available. For fixed \(X\in TG^k\) set \[f_t(y)=\pi_G\operatorname{Fl}_t^S(T R_y X)=L_{\gamma(t)}(y).\] The valid part gives \(C^k\) dependence on \(y\), with \(D_e f_t\) continuous in \(t\). At zero this derivative is invertible. For small \(|t|\) it stays invertible, and right equivariance propagates invertibility to every point. The inverse function theorem and global bijectivity then give a \(C^k\) inverse.\par
Impact and repair boundary. This does not establish the entire flow interval: invertibility holds on an open set of times, but closedness requires further control. Uniform inverse-derivative bounds near finite boundary times would suffice, as would a separate inverse-regularity theorem for the specific group. Neither is established here. The global no-loss and \(G^k\)-completeness proof dependencies remain unresolved. The completeness statements for \(G\) itself have a separate argument and are not refuted by this example.\par
\section*{The Banach model gap and local corrections}
Problem 4 (Major: an invalid boundedness inference in the Banach model construction).\par
Location. Lemma 10.3 (pp. 34–35), its uses in 10.7 and 10.9 (pp. 35–37), and the model and chart constructions in 10.8–10.11 leading to Theorem 3.4 and Lemma 3.7 (pp. 36–39).\par
The proof infers boundedness in \(C^\infty(U,E)\) from pointwise bounded values. The functions \(f_n(x)=\sin(nx)\) on \((-1,1)\) are bounded by one, while \(f_n'(0)=n\). Thus the general function-space implication is false in the compact-open jet topology. This is not a half-Lie-group counterexample to 10.3.\par
Severity and actual dependency. Lemma 10.7 invokes 10.3 both to compare the identity-jet norm with the compact-open topology and to prove completeness. Lemma 10.8 uses that Banach model; 10.9 invokes 10.3 again, 10.10 uses 10.7 to identify the jet model, and 10.11 relies on 10.10. On p. 39, Theorem 3.4 follows explicitly from 10.8, 10.9, and 10.11, while Lemma 3.7 uses the charts of 10.8. The false inference is therefore part of the construction proving the main differentiable-element result. This is a Major proof gap. The new replacement below is counterevidence to irreparability, not evidence that the original argument was valid.\par
Under the local-addition hypothesis already in the main theorem, one can obtain the compact-parameter bounds needed for fixed \(k\). Normalize \(\tau\) and put \(\Omega(y,z)=\tau_y^{-1}(z)\). Then \[L_yv:=T_eR_yv=\Omega(y,\tau_e(v)y),\qquad
 L_y^{-1}w=\Omega(e,\tau_y(w)y^{-1}).\] These are jointly continuous in local tangent trivializations and linear in the tangent argument. Joint continuity at zero implies local operator-norm boundedness; finite covers give bounds on compact parameter sets. The relation \[R_y=\tau_y\circ L_y\circ\tau_e^{-1}\] bounds \(j_e^{k+1}R_y\) and \(j_y^kR_y^{-1}\) for \(y\) in a fixed compact set, using smoothness of \(\tau\). Together with the corrected polynomial jet bounds below, this supplies the estimate used in 10.7 to compare the full identity-jet norm with compact-open jets. Cauchy sequences have continuous derivative limits; integration along small segments identifies those limits as derivatives. Right invariance passes to the limit.\par
Boundary. This verifies replacement bounds under \(\tau\), not the stronger literal statement of 10.3 with one neighbourhood for all \(k\) and all compact sets, without that hypothesis. The norm in 10.7 includes the intermediate derivatives: its printed ellipsis was checked visually. For 10.8, use actual \(C^k\) diffeomorphisms as chart centres and shrink the model neighbourhood so that the identity derivative remains invertible. The inverse function theorem then provides \(C^k\) inverses of the resulting left translations. This avoids assuming an inverse is \(C^k\) for an arbitrary right-invariant \(C^k\) function. The fixed-\(k\) Banach model and these local charts can thus be supported, but the complete Banach/ILB proof is not claimed to have been rewritten.\par
Problem 5 (Minor: normalized jet evaluation and bundle transitions).\par
Location. Definitions 9.1 and 9.4, Lemma 9.5 (pp. 29–31), and Definition 9.11 (p. 34).\par
The normalization is \(p_j=D^jf/j!\). For \(1\le j\le k-1\), the normalized \(j\)th derivative of \(Tf\) is \[\left(p_j(\pi_H^j),\
 (j+1)p_{j+1}(\pi_H^j,X)
 +j p_j(\pi_H^{j-1},\pi_V)\right),\] with symmetrized mixed terms. The printed first term misses \(j+1\), in particular \(k\) in the last component. For \(f(x)=x^2\), \(k=2\), \(x=0\), one has \(Tf(x,X)=(x^2,2xX)\); the horizontal derivative is \((0,2X)\), whereas the printed formula gives \((0,X)\). The old constants in 9.5 cannot simply be retained after correction: at \(X=4\) the actual sum norm is \(12\), while its first displayed bound is \(11\).\par
For the same sum/product norms, sufficient corrected bounds are \[\begin{aligned}
 \|\sigma\odot\xi\|&\le
 \|\xi\|+k\|\sigma\|(1+\|\xi\|),\\
 \|\widetilde\sigma\odot\widetilde\xi-\sigma\odot\xi\|
 &\le k(1+\|\widetilde\xi\|)\|\widetilde\sigma-\sigma\|
 +(1+k\|\sigma\|)\|\widetilde\xi-\xi\|.\end{aligned}\] Collecting each \(p_j\) contribution proves these inequalities. The polynomial-control requirement in 10.6 survives. The original proof on p. 31 already lists the correct unnormalized derivatives, and 10.6 needs a fixed-order polynomial bound rather than the incorrect precise constants. These facts support the Minor scope.\par
Definition 9.11 freezes the base variable in a fibre coordinate change. On a trivial real line bundle take \(\psi(x,y)=y+x\). The zero section has first jet \((0,0,0)\) at zero. After changing coordinates the section is \(x\), with derivative one; freezing \(\psi(0,y)=y\) yields derivative zero. The transition must use the full jet \(j^k\psi\bullet j^k(\operatorname{Id},s)\). Also the right side of Lemma 9.3 must be \(g\circ f\) for \(j^kg\bullet j^kf\). These are local formula corrections, not counterexamples to the intended jet bundles or main group theorem.\par
Problem 6 (Minor: inconsistent bracket signs).\par
Location. Lemma 3.8(a)–(c) and proof (pp. 6–7); the related identity in Lemma 7.4 (p. 23).\par
On the positive affine group \((a,b)(c,d)=(ac,ad+b)\) take \[R_X=a\partial_a+b\partial_b,\qquad R_Y=\partial_b.\] Part (a) gives \(dR_Y\,X-dR_X\,Y=-Y\). But \((\operatorname{Fl}_t^{R_X})^*R_Y=e^{-t}R_Y\), so the negative derivative in (c) gives \(+Y\). Deleting that minus sign makes (c) agree with (a), (b), and the pullback calculation in the proof.\par
A consistent choice is also needed in 7.4: its identity \(-[R_X,R_Y]=R_{[X,Y]}\) uses the opposite convention to 3.8(a). Alternatively, retain the conventional Lie-algebra sign and negate 3.8(a)–(b), adjusting \(\operatorname{ad}\) consistently. This verifies the inconsistency, not an unqualified global sign substitution through all curvature and Eulerian formulas. The existence arguments do not depend on the convention chosen.\par
Problem 7 (Minor: an extension equivalence has the wrong direction).\par
Location. Theorem 5.10(d), final map (p. 12).\par
Let \(G=N=(\mathbb R,+)\), \(\alpha=\operatorname{Id}\), \(f=0\), and \(b(t)=t^2\). The printed transformation gives \(f_1(t,s)=-2ts\). The stated map \(E_f\to E_{f_1}\) is \(\phi(t,n)=(t,n+t^2)\), but \[\phi((1,0)(1,0))=(2,4),\qquad
 \phi(1,0)\mathbin{\cdot_{f_1}}\phi(1,0)=(2,0).\] The general discrepancy is \(4ts\). The correct forward map is \((t,n)\mapsto(t,n-t^2)\); in multiplicative notation it is \((x,n)\mapsto(x,b(x)^{-1}n)\). The printed map has the reverse direction. Substitution verifies this repair; it does not fix Problem 1.\par
Problem 8 (Minor: flatness gives a semidirect product here).\par
Location. Remark 6.3, final sentence (p. 18).\par
Use the trivial principal \(S^1\)-bundle over the compact simply connected manifold \(S^2\), with its flat connection. The homomorphic section is \(s(\varphi)(x,z)=(\varphi(x),z)\). The gauge kernel is the abelian group \(C^\infty(S^2,S^1)\). Conjugation acts by \(h\mapsto h\circ\varphi^{-1}\), generally nontrivially: take \(h(x)=e^{ix_3}\) and a rotation reversing \(x_3\). This distinguished kernel is not central, whereas it would be central in its direct product with \(\operatorname{Diff}(S^2)\). The section supplies a semidirect product. This refutes the direct-product assertion as an equivalence of these extensions, without asserting a stronger abstract non-isomorphism. Replacing direct by semidirect repairs this particular conclusion.\par
\section*{Counterevidence, coverage, and limits}
Theorem 7.7 on \(G\) itself has positive evidence with its exact hypotheses. A strong metric supplies Hilbert charts and local equivalence of metric and coordinate norms. Right translations are isometries, so a Cauchy sequence can be moved into a small fixed normal neighbourhood. Short paths cannot leave a larger normal ball; this supplies the local distance comparison in the proof. Metric completeness supports geodesic completeness.\par
For (e), the endpoint set is weakly closed by hypothesis; its intersection with an energy ball is weakly compact in \(L^2\). Weak lower semicontinuity gives a minimizer without requiring convexity of the endpoint set. Zero control handles equal endpoints. The strong metric also supplies continuous tangent right translation locally through \[T_xR_h(v)=t^{-1}\log_{xh}\bigl((\exp_x(tv))h\bigr)\] for small fixed \(t\). With the energy identity this justifies the \(H^1\)-path/\(L^2\)-control correspondence in the direct method. Local minimization and Gauss’ lemma yield a smooth geodesic. No defect in this argument on \(G\) with the extra conditions was established; the separate \(G^k\) dependency is Problem 3.\par
The initial review covered all 42 pages: definitions, differentiable elements and evolution, extensions and bundle examples, metrics and geodesics, local additions and jets, right-invariant mapping spaces, final proofs, and references. Exact page ranges are retained in coverage. External proofs on Sobolev groups, convenient calculus, and curvature were traced but not re-proved in full. The ILB definition does not require density; nondensity is not a defect. Finite-level uniqueness gives compatibility for the smooth inverse-limit evolution. No obligatory extra spatial derivative loss was established in Theorem 4.2.\par
Two reading artefacts were excluded. The near-minimizer formula on p. 26 has absolute-value bars and is correct; a text-only suggestion of a reversed inequality was withdrawn. The norm on p. 35 has an ellipsis for intermediate derivatives. Both were resolved by actual source-image inspection.\par
\section*{Supporting checks}
The assessment is Major: Problems 1–4 concern two false extension assertions and two core proof gaps; Problems 5–8 are Minor. Problem 4’s fixed-order repair does not lower its original severity. No claim of whole-paper falsity or irreparability is made. Author adoption is unknown.\par
During the initial review, the root executor performed the synthesis, analytic checks, exact symbolic calculations, report writing, compilation, and actual source/report page inspection. Two native helpers supplied bounded checks of Sections 3–5 and 8–10, and Sections 5.11–5.12 and 7–8, respectively. Their corrections and limits are retained. Focused independent acceptance rechecked all eight Problems, their source hypotheses, dependencies, counterevidence, and repair limits. The current executor revised Problem 4, rebuilt and parsed the report, and inspected every revised page. The native read-only challenge superseded its provisional Minor view after tracing the core dependencies; it supplied no compilation, page viewing, or human review.\par
The primary PDF has 502571 bytes and SHA256\par
> `a155b9c7c81026f94ae5368bc687835259eccde3c336683de433994f2d6ca4ac`.\par
The complete extracted text has 2084 lines and SHA256\par
> `599be3b094dbcaa8720f75654f65e4281c1a7d7f0c9c6dc5d4e89fee3e2e08ee`.\par
The exact-check script verifies the local-addition inverse, normalization, cocycle maps, bracket signs, quadratic jet, bundle transition, and spray ODE; its actual exit code is zero. The infinite Hilbert-orbit obstruction is an analytic proof; finite projection norms are illustrative only.\par
The exact TeX, compiler/version, dependency snapshot, before/after facts, real exit codes, logs, PDF, images, and actual page-view receipts are bound by the inventory. The controller’s UI observation is attributed to it. Runtime self-description is separate: the backend was not independently verified and no model equivalence is claimed. Scientific adjudication remains pending; `human\_reviews=[]`.\par
9 M. Bauer, P. Harms, and P. W. Michor, Regularity and completeness of half-Lie groups, arXiv:2302.01631v3, 7 January 2025. https://arxiv.org/abs/2302.01631v3. EMS Press, article record for DOI 10.4171/JEMS/1587. https://ems.press/journals/jems/articles/14298489.\par
\end{document}
